\chapter{Flat morphisms}\label{IV}
% original p. 87
We give here mainly the flatness properties which have served us in
the preceding expos\'es. A more detailed study will be found in
Chapter~IV of the ``\'El\'ements de G\'eom\'etrie Alg\'ebrique'' in
preparation\footnote{\Cf EGA IV 11 and 12.}, where one studies
systematically the following situation: $X$ being locally of finite
type over~$Y$ locally noetherian, and~$\cal{F}$ coherent on~$X$ and
$Y$-flat, to give relations between the properties of~$Y$, those
of~$\cal{F}$, and those of the coherent sheaves induced by~$\cal{F}$
on the fibers of~$X\to Y$ (from the point of view notably of the
dimension, of the cohomological dimension, of the depth etc.). One
has in particular a systematic way of obtaining theorems of
\emph{Seidenberg} or \emph{Bertini} type (for hyperplane sections).
The essential result for the application of flatness methods in this
context is the following (which will be proved below): If~$Y$ is
integral, $X$ of finite type over~$Y$, $\cal{F}$ coherent on~$X$,
there exists a nonempty open set~$U$ of~$Y$ such that~$\cal{F}$ is
$Y$-flat at the points of~$X$ above~$U$. A second way, doubtless even
more important, in which flatness enters Algebraic Geometry, is
\emph{descent theory}: see for example the two expos\'es of
Grothendieck on the subject at the S\'eminaire
Bourbaki\footnote{and, for a more detailed account,
Expos\'es~\Ref{VIII} and~\Ref{IX} below.}. Flatness thus seems to be
one of the central technical notions in Algebraic Geometry.

Let us recall that the notion of flatness and faithful flatness was
introduced by Serre in GAGA. An account of the following
nos.~\Ref{IV.1}~and~\Ref{IV.2} will also be found in Bourbaki's
Alg.\ Comm.\ (which of course, as the title of the book indicates,
does not confine itself to the case of commutative base
rings)\footnote{N. Bourbaki, Alg\`ebre Commutative, Chap\ptbl I
(Modules Plats), Act.\ Sc.\ Ind 1290, Paris, Hermann (1961).}.

Contrary to the preceding expos\'es, we do not assume that the rings
considered are necessarily noetherian.

\section{Sorites on flat modules}
\label{IV.1}

A module $M$ over the ring~$A$ is called \emph{flat} (or $A$-flat if
one wants to specify~$A$)
% original p. 88
if the functor
$$
T_M \colon N \mto M \otimes_{A}N
$$
(which is in any case right exact) is \emph{exact}, \ie transforms
monomorphisms into monomorphisms. It amounts to the same to say that
the first right derived functor, or all the right derived functors,
are zero, \ie that one has
$$
\Tor^{A}_{1}(M,N)=0 \qquad \text{for all $N$,}
$$
\resp that one has
$$
\Tor^{A}_{i}(M,N)=0 \qquad \text{for $i>0$, all $N$.}
$$
Since the $\Tor_{i}$ commute with inductive limits, it suffices
moreover to verify these conditions for $N$ of finite type, and even
(taking then a composition series of~$N$ with monogenous quotients)
that one has
$$
\Tor^{A}_{1}(M,N)=0
$$
if $N$ is monogenous, \ie of the form $A/I$, $I$ being an ideal
of~$A$. Let us note moreover that
$$
\Tor^{A}_{1}(M,A/I)=0 \quad\Longleftrightarrow\quad I\otimes_{A}M \rightarrow
M=A\otimes_{A}M\text{ is \emph{injective}},
$$
as one sees from the exact sequence of the $\Tor$, taking into account
$\Tor^{A}_{1}(M,A)=0$. Hence $M$ flat is equivalent to saying that for
every ideal~$I$, the natural homomorphism
$$
I\otimes_{A}M \to IM
$$
is an isomorphism. It suffices moreover to verify this for $I$ of
finite type; a fortiori it suffices to verify that the functor
$M\otimes$ is exact on modules \emph{of finite type}.

As whenever one has an exact functor $T$, if one identifies, for a
subobject $N'$ of~$N$, $T(N')$ with a subobject of~$T(N)$, one has
$$
T(N' \cap N'')=T(N') \cap T(N'')
$$
$$
T(N' + N'')=T(N') + T(N'')
$$
for two subobjects $N'$, $N''$ of~$N$.

A direct sum of flat modules, a direct factor of a flat module, is
flat. In particular, $A$ being flat, a \emph{free} module, hence also
a \emph{projective} module, is flat. The tensor product of two flat
modules is flat, and if $M$ is flat over~$A$, then $M\otimes_{A}B$ is
flat over~$B$ for every change of base $A \to B$ (because of the
associativity of the tensor product and of the fact that a composite
of
% original p. 89
exact functors is exact). If $M$ is flat over~$B$, $B$ flat over~$A$,
then $M$ is flat over~$A$ (same reason).

The exact sequence of the $\Tor$, plus the ``commutativity'' of the
$\Tor$, gives:

\begin{proposition}
\label{IV.1.1}
Let $0 \to M' \to M \to M'' \to 0$ be an exact sequence of~$A$-modules,
$M''$ being flat. Then
\begin{enumerate}
\item[(i)] this sequence remains exact upon tensoring with any
$A$-module $N$
\item[(ii)] for $M$ to be flat, it is necessary and sufficient that
$M'$ be so.
\end{enumerate}
\end{proposition}

One can therefore say that from the point of view of behavior under
tensor products, flat modules are ``as good'' as free or projective
modules (and the exact sequence of~\Ref{IV.1.1} in particular is ``as
good'' as if it split).

Let $S$ be a multiplicatively stable subset of~$A$; then $S^{-1}A$ is
flat over~$A$, since $S^{-1}A \otimes N=S^{-1}N$ is an exact functor
in~$N$. If $M$ is $A$-flat, then $S^{-1}M=S^{-1}A \otimes M$ is
$S^{-1}A$-flat, the converse being true if $M \to S^{-1}M$ is an
isomorphism, \ie if the $s \in S$ are bijective in $M$ (because of the
transitivity of flatness, $S^{-1}A$ being flat over~$A$). More
generally, the case of a morphism of preschemes $X \to Y$ and of a
quasi-coherent sheaf $\cal{F}$ on~$X$ whose flatness with respect
to~$Y$ one wants to study leads to the situation with two rings:

\begin{proposition}
\label{IV.1.2}
Let $A \to B$ be a homomorphism of rings, $M$ a $B$-module, $T$ a
multiplicatively stable subset of~$B$.
\begin{enumerate}
\item[(i)] If $M$ is $A$-flat, then $T^{-1}M$ is $A$-flat (hence also
$S^{-1}A$-flat for every multiplicatively stable subset $S$ of~$A$
mapping into $T$).
\item[(ii)] Conversely, if $M_{\goth{n}}$ is flat over~$A_{\goth{n}}$
for every maximal ideal $\goth{n}$ of~$B$, $M_{\goth{n}}$ is flat
over~$A$ (or, what amounts to the same, over~$A_{\goth{m}}$, where
$\goth{m}$ is the prime ideal of~$A$ inverse image of~$\goth{n}$)
then $M$ is $A$-flat.
\end{enumerate}
\end{proposition}

Indeed one has the formula, functorial with respect to the
$A$-module~$N$:
$$
T^{-1}M \otimes_{A} N=T^{-1}(M \otimes_{A} N)
$$
since the two sides are functorially isomorphic to $T^{-1}B
\otimes_{B} M \otimes_{B} N_{(B)}$ (with $N_{(B)}=N \otimes_{A} B$) by
virtue of the associativity formulas of~$\otimes$. It follows at once
that if $M \otimes_{A} N$ is exact in~$N$, the same is true
of~$T^{-1}M \otimes_{A} N$ (as a composite of two exact functors),
whence~(i). And (ii) follows likewise, since to verify the exactness
of a sequence of~$B$-modules, it suffices to verify the exactness of
the localizations at all the maximal ideals of~$B$.

\begin{proposition}
\label{IV.1.3}
\begin{enumerate}
\item[(i)] Let
% original p. 90
$M$ be a flat $A$-module. If $x \in A$ is not a divisor of~$0$ in $A$,
it is not a divisor of~$0$ in $M$. In particular, if $A$ is integral,
$M$ is torsion-free.
\item[(ii)] Suppose that $A$ is integral and that for every maximal
ideal $\goth{m}$ of~$A$, $A_{\goth{m}}$ is principal (for example $A$
a Dedekind ring, or even principal). For the $A$-module $M$ to be
flat, it is necessary and sufficient that it be torsion-free.
\end{enumerate}
\end{proposition}

One obtains (i) by noting that the homothety $x$ in $M$ is obtained by
tensoring with~$M$ the homothety $x$ in $A$. For (ii), one may already
assume $A$ principal thanks to~\Ref{IV.1.2} (ii); one must show that
if $M$ is torsion-free, then for every ideal $I$ of~$A$, the injection
$I \to A$ tensored with $M$ is an injection, which means that the
generator $x$ of~$I$ is not a divisor of~$0$ in~$M$, O.K.

\section{Faithfully flat modules}
\label{IV.2}

A functor $F$ from one category into another is called
\emph{faithful} if for all $X$, $Y$, the map $\Hom (X,Y) \to \Hom
(F(X),F(Y))$ is injective. If it is an additive functor of additive
categories, it amounts to the same to say that $F(u)=0$ implies
$u=0$, and this implies that $F(X)=0$ implies $X=0$. For $F$ to be
\emph{faithful and exact}, it is necessary and sufficient that the
following condition be satisfied: for every sequence $M' \to M \to
M''$ of morphisms in $\cal{C}$, the transformed sequence
$F(M') \to F(M) \to F(M'')$
is exact \emph{if} and \emph{only if} the preceding one is. Or again:
$F$ is exact, and $F(X)=0$ implies $X=0$. (N.B. to be able to speak
of exactness, one must assume the categories in question
\emph{abelian}). Suppose one has a family $(M_{i})$ of nonzero objects
of~$\cal{C}$ such that every nonzero object of~$\cal{C}$ has a
subobject admitting a quotient isomorphic to one of the $M_{i}$. Then
$F$ faithful and exact is equivalent to: $F$ exact, and $F(M_{i}) \neq
0$ for all~$i$. If $\cal{C}$ is the category of modules over a
ring~$A$, one can take for example for $(M_{i})$ the family of the
$A/\goth{m}$, $\goth{m}$ running through the maximal ideals of~$A$.
(Indeed, every nonzero module admits a nonzero monogenous submodule,
hence isomorphic to some $A/I$, $I$ an ideal~$\neq A$, which by Krull
admits a quotient $A/\goth{m}$). From these sorites, one deduces in
particular:

\begin{proposition}
\label{IV.2.1}
Let $M$ be an $A$-module. Equivalent conditions:
\begin{enumerate}
\item[(i)] The functor $M \otimes_{A}$ is faithful and exact.
\item[(i~bis)]
% original p. 91
$M$ is flat, and $M \otimes_{A} N=0$ implies $N=0$
\item[(i~ter)] $M$ is flat, and $M \otimes A/\goth{m} \neq 0$ for
every maximal ideal $\goth{m}$ of~$A$.
\item[(ii)] For every sequence of homomorphisms $N' \to N \to N''$,
the sequence tensored with $M$ is exact if and only if the initial
sequence is.
\end{enumerate}
\end{proposition}

One then says that $M$ is a \emph{faithfully flat} $A$-module. In
particular, if $M$ is faithfully flat, then $N \to N'$ is a
monomorphism (epimorphism, isomorphism) if and only if the
homomorphism tensored with~$M$ is. A faithfully flat module is
\emph{faithful}, since the homothety $f$ in~$M$ is obtained by
tensoring with~$M$ the homothety $f$ in~$A$.

One sees as in~\Ref{IV.1} the usual transitivity properties: the
tensor product of two faithfully flat modules is faithfully flat; if
$M$ is faithfully flat over~$A$, $M \otimes_{A}B$ is faithfully flat
over~$B$ for every extension of the base $A \to B$; if $B$ is an
$A$-algebra which is faithfully flat over~$A$ and if $M$ is a
faithfully flat $B$-module, it is a faithfully flat $A$-module.

\begin{corollary}
\label{IV.2.2}
Let $A \to B$ be a local homomorphism of local rings, $M$ a
$B$-module of finite type. For $M$ to be faithfully flat over~$A$, it
is necessary and sufficient that it be flat over~$A$ and nonzero.
\end{corollary}

This follows from criterion (i~ter) and from Nakayama. In particular,
\emph{for $B$ to be $A$-flat, it is necessary and sufficient that it
be faithfully $A$-flat}.

\begin{proposition}
\label{IV.2.3}
Let $A \to B$ be a homomorphism of rings, $M$ a $B$-module which is
faithfully flat over~$A$. For every prime ideal $\goth{p}$ of~$A$,
there exists a prime ideal~$\goth{q}$ of~$B$ which induces it.
\end{proposition}

Dividing by $\goth{p}$, one is reduced to the case where
$\goth{p}=0$. Localizing at the prime ideal $0$, one is reduced to the
case where $A$ is a field. But $M$, being faithfully flat over~$A$, is
nonzero, a fortiori $B \neq 0$, hence $B$ has a prime ideal, which can
only induce the unique prime ideal of~$A$! Geometrically, one can say
that the existence of a quasi-coherent sheaf $\cal{F}$ on~$X=\Spec
(B)$ which is ``faithfully flat'' relative to~$A$ implies that $X \to
Y=\Spec (A)$ is \emph{surjective}.

\begin{corollary}
\label{IV.2.4}
Suppose $M$ flat over~$A$, of finite type over~$B$ and
$\Supp M=\Spec(B)$ (\ie $M_{\goth{q}}\neq 0$ for every prime ideal
$\goth{q}$ of~$B$). Then the prime ideals $\goth{q}$ of~$B$ minimal
among those containing $\goth{p}B$ induce~$\goth{p}$.
\end{corollary}

One
% original p. 92
is again reduced to the case $\goth{p}=0$ (since the hypotheses are
all preserved on dividing), hence $A$ integral. One is reduced to the
following statement:

\begin{corollary}
\label{IV.2.5}
$M$ being as above, every minimal prime ideal $\goth{q}$ of~$B$
induces a prime ideal $\goth{p}$ of~$A$ which is minimal.
\end{corollary}

Indeed, localizing at $\goth{p}$ and $\goth{q}$, one is reduced to
proving that if $A$ and $B$ are local and the homomorphism $A \to B$
local, $M$ a nonzero $B$-module flat over~$A$, and if $B$ is of
dimension $0$, then $A$ is of dimension $0$. By~\Ref{IV.2.2}
and~\Ref{IV.2.3}, one concludes that every prime ideal of~$A$ is
induced by a prime ideal of~$B$, hence by the maximal ideal of~$B$,
hence is the maximal ideal, qed. Geometrically, \Ref{IV.2.5} means
that every irreducible component of $X=\Spec (B)$ dominates some
irreducible component of~$Y=\Spec (A)$ (granted the existence of a
quasi-coherent sheaf of finite type on~$X$, with support~$X$, and flat
with respect to~$Y$).

It will be noted that one did not have to assume in~\Ref{IV.2.4} $M$
faithfully flat over~$A$, but nothing then guarantees the existence of
a prime ideal containing $\goth{p}B$ (hence of a minimal one among
such).

\begin{proposition}
\label{IV.2.6}
Let $i \colon A \to B$ be a homomorphism of rings. Equivalent
conditions:
\begin{enumerate}\item[(i)] $B$ is a faithfully flat $A$-module.
\item[(ii)] $B$ is flat over~$A$, and $\Spec (B) \to \Spec (A)$ is
surjective
\item[(ii~bis)] $B$ is flat over~$A$, and every maximal ideal is
induced by an ideal of~$B$.
\item[(iii)] $i$ is injective and $\Coker i$ is a flat $A$-module.
\item[(iv)] The functor $M_{(B)}=M \otimes_{A} B$ in the $A$-module
$M$ is exact, and the canonical functorial homomorphism $M \to
M_{(B)}$ is injective.
\item[(iv~bis)] For every ideal $I$ of~$A$, $I \otimes_{A} B \to IB$
is an isomorphism, and the inverse image of~$IB$ in $A$ is equal
to~$I$.
\end{enumerate}
\end{proposition}

One has (i)$\To$(ii) by~\Ref{IV.2.3}, (ii)$\To$(ii~bis) is trivial,
(ii~bis)$\To$(i) by criterion (i~ter) of~\Ref{IV.2.1}. One has
(iii)$\To$(iv) by~\Ref{IV.1.1}, (iv)$\To$(iv~bis) trivially (by
taking $M=A/I$ in the second condition~(iv~bis)), and
(iv~bis)$\To$(i) by virtue of the criterion of flatness by ideals seen
at the beginning of~\Ref{IV.1} and of criterion \Ref{IV.2.1}~(i~ter).
Finally, (iv)$\To$(iii) by an easy converse of~\Ref{IV.1.1}, and
(i)$\To$(iv) since
% original p. 93
if $N$ is the kernel of~$M \to M \otimes_{A}B =T(M)$, then ($T$~being
exact) $N \to T(N)$ is zero, whence $T(N)=N \otimes_{A}B =0$, whence
$N=0$, qed.

\section{Relations with completion}
\label{IV.3}
Let $A$ be a noetherian ring, $I$ an ideal in $A$, $\widehat{A}$ the
separated completion of~$A$ for the $I$-preadic topology, and for
every $A$-module $M$, let $\widehat{M}$ be its completion for the
$I$-preadic topology. It is an $\widehat{A}$-module, whence a
canonical homomorphism
$$
M \otimes_A \widehat{A} \to \widehat{M}.
$$
When $M$ runs through the modules \emph{of finite type}, the functor
$M \mto \widehat{M}$ is exact, as follows easily from \emph{Krull's
theorem: If $N \subset M$, the topology of~$N$ is the one induced by
the topology of~$M$}. Since $M \otimes_A \widehat{A}$ is right exact,
one easily concludes from this (by resolving $M$ by $L \to L' \to M$,
with $L$ and $L'$ free of finite type) that the functorial
homomorphism above is an \emph{isomorphism} ($\widehat{M}$ being also
right exact) and consequently that $M \otimes_A \widehat{A}$ is also
an \emph{exact} functor in $M$. Hence:
\begin{proposition}
\label{IV.3.1}
Let $A$ be a noetherian ring, $I$ an ideal of~$A$; then the separated
completion $\widehat{A}$ of~$A$ (for the $I$-preadic topology) is
\emph{flat} over~$A$.
\end{proposition}

\begin{corollary}
\label{IV.3.2}
For $\widehat{A}$ to be faithfully flat over~$A$, it is necessary and
sufficient that $I$ be contained in the radical of~$A$.
\end{corollary}

Indeed, it suffices to apply criterion \Ref{IV.2.1} (i~ter).

These results summarize all that one knows how to say, from the point
of view of linear algebra, about the relations between $A$ and
$\widehat{A}$. Corollary \Ref{IV.3.2} is used above all when $A$ is a
noetherian local ring and $I$ is contained in the maximal ideal
$\goth{m}$ (and most often, is equal to it).

\section{Relations with free modules}
\label{IV.4}

\begin{proposition}
\label{IV.4.1}
Let $A$ be a ring, $I$ an ideal of~$A$, $M$ an $A$-module. Suppose
that one is under one or the other of the following hypotheses:
\begin{enumerate}
\item[(a)] $I$ is nilpotent
\item[(b)] $A$ is noetherian, $I$ is in the radical of~$A$, and $M$ is
of finite type.
\end{enumerate}
For $M$ to be free over~$A$, it is necessary and sufficient that $M
\otimes A/I$ be free over~$A/I$, and that $\Tor_1^A(M,A/I) = 0$.
\end{proposition}

It is
% original p. 94
necessary; let us prove sufficiency. Let $(e_i)$ be a family of
elements of~$M$ whose image in $M \otimes A/I = M/IM$ defines a basis
of it over~$A/I$ (it is a finite family in case (b)). Let $L$ be the
free $A$-module constructed on the same index set; one thus has a
homomorphism $L \to M$ such that tensoring $T$ with $A/I$ induces an
isomorphism $T(L) \isomto T(M)$. If $Q$ is the cokernel of~$L \to M$,
one therefore has $T(Q) = 0$, whence $Q=0$ by virtue of Nakayama
(valid under one or the other condition (a) or (b)). Hence $L \to M$
is surjective; let $R$ be its kernel; one thus has an exact sequence
$$
0 \to R \to L \to M \to 0
$$
whence, since $\Tor_1^A(M,A/I)=0$, an exact sequence $0 \to T(R) \to
T(L) \to T(M) \to 0$, whence $T(R)=0$, whence further $R=0$ by virtue
of Nakayama (taking into account that in case (b), $R$ is of finite
type since $A$ was assumed noetherian).

\begin{corollary}\label{IV.4.2}
One can replace the condition $\Tor_1^A(M,A/I)=0$ by: the canonical
surjective homomorphism
\begin{equation*}
\label{eq:IV.2.*}
\tag{$*$} { \gr_I^0(M) \otimes_{A/I} \gr_I(A) \to \gr_I(M) }
\end{equation*}
is an isomorphism.
\end{corollary}

Indeed, if $M$ is free, this is certainly satisfied. One must
therefore prove that if $M \otimes A/I$ is free over~$A/I$ and the
condition on the $\gr$ is satisfied, then $M$ is free. One takes up
the proof above again, constructing $L \to M$; it follows from the
hypothesis that this homomorphism induces an isomorphism for the
associated graded modules, hence its kernel is contained in the
intersection of the $I^nL$, hence is zero (as is trivial in (a), and
well known in (b)). Qed.

\begin{corollary}
\label{IV.4.3}
Suppose that $A/I$ is a field. Then the following conditions on~$M$
are equivalent:
\begin{enumerate}
\item[(i)] $M$ is free
\item[(ii)] $M$ is projective
\item[(iii)] $M$ is flat
\item[(iv)] $\Tor_1^A(M, A/I) = 0$
\item[(v)] The canonical homomorphism \eqref{eq:IV.2.*} is bijective.
\end{enumerate}
\end{corollary}

Indeed, in the case considered, $M \otimes A/I$ is automatically
free.

The preceding result is valid in the following two cases:
\begin{enumerate}
\item[(a)] $M$ is an \emph{arbitrary} module over a local ring $A$
whose maximal ideal $I$ is \emph{nilpotent} (for example an artinian
local ring).
\item[(b)] $M$ is a module \emph{of finite type} over a
\emph{noetherian local} ring.
\end{enumerate}

Let us recall,
% original p. 95
for the record:
\begin{corollary}
\label{IV.4.4}
Suppose that $A$ is a \emph{noetherian local integral} ring with
maximal ideal $\goth{m} = I$, residue field $k = A/I$, field of
fractions $K$. Let $M$ be a module of finite type over~$A$. Then the
equivalent conditions \textup{(i)} to \textup{(v)} above are also
equivalent to
\begin{enumerate}
\item[(vi)] $M \otimes_A K$ and $M \otimes_A k$ are vector spaces of
the same dimension (\ie the rank of~$M$ over~$A$ is equal to the
minimum number of generators of the $A$-module $M$).
\end{enumerate}
\end{corollary}

Immediate proof; we leave to the reader the task of generalizing to
the case where~$A$ is only assumed to be without nilpotent elements:
one must then require that the ranks of~$M$ at the minimal prime
ideals of~$A$ be equal to the dimension of the vector space $M
\otimes_A k$.
